

%% %%%%%%%%%%%%%%%%%%%%%%%%%%%%%%%%%%%%%%%%%%%%%%%%%
%% Template for a conference paper, prepared for the
%% Food and Resource Economics Department - IFAS
%% UNIVERSITY OF FLORIDA
%% %%%%%%%%%%%%%%%%%%%%%%%%%%%%%%%%%%%%%%%%%%%%%%%%%
%% Version 1.0 // November 2019
%% %%%%%%%%%%%%%%%%%%%%%%%%%%%%%%%%%%%%%%%%%%%%%%%%%
%% Ariel Soto-Caro
%%  - asotocaro@ufl.edu
%%  - arielsotocaro@gmail.com
%% %%%%%%%%%%%%%%%%%%%%%%%%%%%%%%%%%%%%%%%%%%%%%%%%%
\documentclass[12pt]{article}
\usepackage{UF_FRED_paper_style}
\usepackage{threeparttable}
\usepackage{booktabs}
\usepackage{tabularx}
\usepackage{amsmath} 
\usepackage{lipsum}  %% Package to create dummy text (comment or erase before start)
\usepackage{graphicx}

\usepackage{amsthm}
\usepackage[dvipsnames]{xcolor}
\theoremstyle{plain}
\newtheorem{thm}{Theorem}
\newtheorem{defi}{Definition}
\newtheorem{ass}{Assumption}
\newtheorem{prop}{Proposition}
\newtheorem{lem}{Lemma}
\newtheorem{coro}{Corollary}

\usepackage{hyperref}					% Reference
\usepackage[nameinlink]{cleveref} % nameinlink optional but nice
\hypersetup{colorlinks,%				% Reference color setup	
	linkcolor=red,%
	citecolor=blue}
% Tell cleveref how to name this environment
\crefname{thm}{Theorem}{Theorems}
\Crefname{thm}{Theorem}{Theorems}
\crefname{ass}{Assumption}{Assumptions}
\Crefname{ass}{Assumption}{Assumptions}
\crefname{defi}{Definition}{Definitions}
\Crefname{defi}{Definition}{Definitions}
\crefname{prop}{Proposition}{Propositions}
\Crefname{prop}{Proposition}{Propositions}
\crefname{lem}{Lemma}{Lemmas}
\Crefname{lem}{Lemma}{Lemmas}
\crefname{coro}{Corollary}{Corollaries}
\Crefname{coro}{Corollary}{Corollaries}

\newcommand{\red}[1]{{\color{red} #1}}
\newcommand{\blue}[1]{{\color{blue} #1}}
\newcommand{\orange}[1]{{\color{orange} #1}}
\newcommand{\gray}[1]{{\color{gray} #1}}
\newcommand{\green}[1]{{\color{OliveGreen} #1}}
\newcommand{\proptoinverse}{\mathrel{\mskip1mu\reflectbox{$\propto$}\mskip-1mu}}
\usepackage{dcolumn}
\newcolumntype{d}[1]{D..{#1}}
\newcommand{\mc}[1]{\multicolumn{1}{c@{}}{#1}}
\newcommand{\mX}[1]{\multicolumn{1}{X}{\centering #1}}
%% ===============================================
%% Setting the line spacing (3 options: only pick one)
% \doublespacing
% \singlespacing
\onehalfspacing
%% ===============================================

\setlength{\droptitle}{-5em} %% Don't touch

% %%%%%%%%%%%%%%%%%%%%%%%%%%%%%%%%%%%%%%%%%%%%%%%%%%%%%%%%%%
% SET THE TITLE
% %%%%%%%%%%%%%%%%%%%%%%%%%%%%%%%%%%%%%%%%%%%%%%%%%%%%%%%%%%

% TITLE:
\title{Two-country Version of \cite{hirano2025technological}
%\thanks{}
}

% AUTHORS:
\author{
    }
    
% DATE:
\date{\today}

% %%%%%%%%%%%%%%%%%%%%%%%%%%%%%%%%%%%%%%%%%%%%%%%%%%%%%%%%%%
% %%%%%%%%%%%%%%%%%%%%%%%%%%%%%%%%%%%%%%%%%%%%%%%%%%%%%%%%%%
\begin{document}
% %%%%%%%%%%%%%%%%%%%%%%%%%%%%%%%%%%%%%%%%%%%%%%%%%%%%%%%%%%
% %%%%%%%%%%%%%%%%%%%%%%%%%%%%%%%%%%%%%%%%%%%%%%%%%%%%%%%%%%
% ABSTRACT
% %%%%%%%%%%%%%%%%%%%%%%%%%%%%%%%%%%%%%%%%%%%%%%%%%%%%%%%%%%
% %%%%%%%%%%%%%%%%%%%%%%%%%%%%%%%%%%%%%%%%%%%%%%%%%%%%%%%%%%
{\setstretch{.8}
\maketitle
}
\section{Model setup}

\subsection{Representative household}

%--------------------------------------------------------
% US household (equities + US bond only; bond-share cost)
%--------------------------------------------------------
\subsubsection{US}

Let
\[
S_t \equiv Q_{US,t}n_{US,t}+Q_{W,t}n_{W,t},
\qquad
\omega_t \equiv \frac{Q_{US,t}n_{US,t}}{S_t},
\qquad
q_{US,t}\equiv \frac{1}{R_{f,t}}
\]
denote, respectively, the market value of equity purchases, the portfolio weight on US equity within equity
holdings, and the date-$t$ price of the one-period US riskless bond (gross risk-free return $R_{f,t}$).
The US household trades US and RoW equities and the US bond only (it does \emph{not} trade the RoW bond).
It chooses
\[
\{C_{y,t},\,n_{US,t},\,n_{W,t},\,B_{US,t+1}\}
\]
to solve
\[
\max
\ (1-\beta)\log C_{y,t}
+\frac{\beta}{1-\gamma}\log E_t\!\left[(C_{o,t+1})^{1-\gamma}\right]
-\frac{\kappa}{2}(\omega_t-\bar\omega)^2
-\eta\,\Upsilon(\theta_t)
\]
subject to
\begin{align*}
C_{y,t}
&+Q_{US,t}n_{US,t}+Q_{W,t}n_{W,t}
+q_{US,t}B_{US,t+1}
= e_{US,t},\\
C_{o,t+1}
&=
(Q_{US,t+1}+D_{US,t+1})n_{US,t}
+(Q_{W,t+1}+D_{W,t+1})n_{W,t}
+B_{US,t+1}.
\end{align*}
Here $\bar\omega\in[1/2,1]$ is the targeted portfolio weight on US equity (home bias). The term
$\eta\,\Upsilon(\theta_t)$ captures a convex cost in the present-value bond share $\theta_t$,
disciplines large bond positions, and preserves the saving separation. We assume that issuance-cost function $\Upsilon$
 satisfies standard regularity and Inada-type conditions.

\begin{ass}[Issuance costs]\label{ass_issuance_costs}
The issuance-cost function $\Upsilon:\,(-\infty,1)\to\mathbb R$ is $C^1$, strictly convex, and satisfies
\[
\lim_{\theta\uparrow 1}\Upsilon'(\theta)=+\infty,
\qquad
\lim_{\theta\downarrow -\infty}\Upsilon'(\theta)=-\infty.
\]
\end{ass}


\begin{proof}
By \eqref{US_theta_FOC_bond_cost},
\[
\eta\,\Upsilon'(\theta_t)
=\beta\,E_t\!\left[\mathcal M_{t,t+1}\bigl(R_{f,t}-R_{p,t+1}\bigr)\right].
\]
Hence $\Upsilon'(\theta_t)\in[-\overline\xi/\eta,\overline\xi/\eta]$. Since $\Upsilon$ is strictly convex,
$\Upsilon'$ is strictly increasing and continuous on $(-\infty,1)$, and by \Cref{ass_issuance_costs} it is onto $\mathbb R$.
Therefore there exist unique $\underline\theta,\overline\theta\in(-\infty,1)$ such that
$\Upsilon'(\underline\theta)=-\overline\xi/\eta$ and $\Upsilon'(\overline\theta)=\overline\xi/\eta$.
Monotonicity of $\Upsilon'$ yields $\theta_t\in[\underline\theta,\overline\theta]$ for all $t\le\tau$, implying
$\theta_t\ge \underline\theta$ and $1-\theta_t\ge 1-\overline\theta>0$.
\end{proof}

\paragraph{Reformulation in terms of total saving, bond share, and equity share.}
Define the present value of the US bond position and total saving:
\[
b_t\equiv q_{US,t}B_{US,t+1},
\qquad
A_t \equiv S_t+b_t,
\qquad
C_{y,t}=e_{US,t}-A_t.
\]
Define the (present-value) bond share in total saving
\[
\theta_t \equiv \frac{b_t}{A_t},
\qquad
S_t=(1-\theta_t)A_t,
\]
and the gross return on total saving
\[
R_{A,t+1}\equiv (1-\theta_t)R_{p,t+1}+\theta_t R_{f,t},
\qquad
R_{f,t}\equiv \frac{1}{q_{US,t}},
\]
where
\[
R_{p,t+1}=\omega_t R_{US,t+1}+(1-\omega_t)R_{W,t+1},
\qquad
R_{i,t+1}=\frac{Q_{i,t+1}+D_{i,t+1}}{Q_{i,t}},
\ \ i\in\{US,W\}.
\]
Old-age consumption satisfies $C_{o,t+1}=A_tR_{A,t+1}$, so the problem can be written as
\[
\max_{A_t,\omega_t,\theta_t}
\ (1-\beta)\log(e_{US,t}-A_t)
+\frac{\beta}{1-\gamma}\log E_t\!\left[(A_tR_{A,t+1})^{1-\gamma}\right]
-\frac{\kappa}{2}(\omega_t-\bar\omega)^2
-\eta\,\Upsilon(\theta_t).
\]
Under log utility in youth and the risk-sensitive aggregator in old age, the saving choice separates and yields
\begin{align}\label{US_saving_rule_bond_cost}
\frac{1-\beta}{e_{US,t}-A_t}=\frac{\beta}{A_t}
\quad\Rightarrow\quad
A_t=\beta e_{US,t},
\qquad
C_{y,t}=(1-\beta)e_{US,t}.
\end{align}

\paragraph{Portfolio FOCs and stock-pricing implications.}
Define the normalized \emph{undistorted} portfolio kernel
\begin{equation}\label{US_K_def}
\mathcal K_{t,t+1}
\equiv
\frac{R_{A,t+1}^{-\gamma}}{E_t\!\left[R_{A,t+1}^{1-\gamma}\right]},
\qquad
E_t\!\left[\mathcal K_{t,t+1}R_{A,t+1}\right]=1.
\end{equation}
The FOC for $\omega_t$ implies
\begin{align}\label{US_omega_FOC_bond_cost}
\beta(1-\theta_t)\,
E_t\!\left[\mathcal K_{t,t+1}\bigl(R_{US,t+1}-R_{W,t+1}\bigr)\right]
=
\kappa(\omega_t-\bar\omega).
\end{align}
The FOC for $\theta_t$ implies
\begin{align}\label{US_theta_FOC_bond_cost}
\beta\,
E_t\!\left[\mathcal K_{t,t+1}\bigl(R_{f,t}-R_{p,t+1}\bigr)\right]
=
\eta\,\Upsilon'(\theta_t).
\end{align}
The following lemma shows how \Cref{ass_issuance_costs} guarantees that $\theta_t$ is bounded.
\begin{lem}[Endogenous bounds on the bond share]\label{lem2}
Suppose \Cref{ass_issuance_costs} holds and the bond-share FOC \eqref{US_theta_FOC_bond_cost} holds with $\eta>0$ and finite LHS,
i.e.,
\[
\left|\,\beta\,E_t\!\left[\mathcal M_{t,t+1}\bigl(R_{f,t}-R_{p,t+1}\bigr)\right]\right|\le \overline\xi.
\]
Then there exist constants $\underline\theta<\overline\theta<1$ such that
\[
\theta_t\in[\underline\theta,\overline\theta]
\qquad\text{for all }t\le\tau.
\]
In particular, $\theta_t$ is bounded below ($\theta_t$ cannot diverge to $-\infty$) and the equity share is
bounded away from zero:
\[
\theta_t \ge \underline\theta>-\infty,
\qquad
1-\theta_t \ge 1-\overline\theta>0,
\qquad t\le\tau.
\]
\end{lem}


For later use, define
\[
\phi_t
\equiv
\frac{\kappa}{\beta(1-\theta_t)}(\omega_t-\bar\omega)
=
\frac{\kappa C_{y,t}}{(1-\beta)S_t}(\omega_t-\bar\omega),
\qquad
\lambda_t
\equiv
\frac{\eta\theta_t}{\beta}\,\Upsilon'(\theta_t).
\]
Using
\[
R_{A,t+1}=(1-\theta_t)R_{p,t+1}+\theta_t R_{f,t}
\]
and \eqref{US_K_def}, the bond-share FOC implies
\begin{align}\label{US_Rp_pricing}
E_t\!\left[\mathcal K_{t,t+1}R_{p,t+1}\right]
=
1-\lambda_t.
\end{align}
Combining \eqref{US_Rp_pricing} with \eqref{US_omega_FOC_bond_cost} yields
\begin{align}
E_t\!\left[\mathcal K_{t,t+1}R_{US,t+1}\right]
&=
1-\lambda_t+\phi_t(1-\omega_t),\label{US_RUS_pricing}\\
E_t\!\left[\mathcal K_{t,t+1}R_{W,t+1}\right]
&=
1-\lambda_t-\phi_t\omega_t.\label{US_RW_pricing}
\end{align}

Therefore the stock-pricing equations can be written as
\begin{align}
Q_{US,t}
&=
E_t\!\left[\widetilde M^{\,US}_{t,t+1}\bigl(Q_{US,t+1}+D_{US,t+1}\bigr)\right],\label{US_AP_bond_cost}\\
Q_{W,t}
&=
E_t\!\left[\widetilde M^{\,W}_{t,t+1}\bigl(Q_{W,t+1}+D_{W,t+1}\bigr)\right],\notag
\end{align}
with
\begin{align*}
\widetilde M^{\,US}_{t,t+1}
&=
\frac{\mathcal K_{t,t+1}}
{1-\lambda_t+\phi_t(1-\omega_t)},\\
\widetilde M^{\,W}_{t,t+1}
&=
\frac{\mathcal K_{t,t+1}}
{1-\lambda_t-\phi_t\omega_t}.
\end{align*}
Thus the kernel that prices the US stock differs from the undistorted kernel $\mathcal K_{t,t+1}$ because of
both the home-bias wedge $\phi_t$ and the bond-share wedge $\lambda_t$.

%--------------------------------------------------------
% RoW household (equities + US bond + RoW bond; convenience yield; NO bond-share costs)
%--------------------------------------------------------
\subsubsection{Rest-of-World (RoW)}

Let
\[
S_t^* \equiv Q_{US,t}n^*_{US,t}+Q_{W,t}n^*_{W,t},
\qquad
\omega_t^* \equiv \frac{Q_{US,t}n^*_{US,t}}{S_t^*},
\qquad
q_{W,t}\equiv \frac{1}{R_{f,t}^W}
\]
where $q_{W,t}$ is the date-$t$ price of the one-period RoW riskless bond (gross return $R_{f,t}^W$).
The RoW household trades both equities, the US bond, and the RoW bond, and chooses
\[
\{C^*_{y,t},\,n^*_{US,t},\,n^*_{W,t},\,B^*_{US,t+1},\,B^*_{W,t+1}\}
\]
to solve
\begin{align*}
\max\quad U_t^*
&=
(1-\beta)\log C^*_{y,t}
+\frac{\beta}{1-\gamma}\log E_t\!\left[\left(C^*_{o,t+1}\right)^{1-\gamma}\right]
-\frac{\kappa}{2}\left(\omega_t^*-\bar\omega^*\right)^2
+\chi\log\!\left(q_{US,t}B^*_{US,t+1}\right),
\end{align*}
subject to
\begin{align*}
C^*_{y,t}
&+Q_{US,t}n^*_{US,t}+Q_{W,t}n^*_{W,t}
+q_{US,t}B^*_{US,t+1}+q_{W,t}B^*_{W,t+1}
= e_{W,t},\\
C^*_{o,t+1}
&=
(Q_{US,t+1}+D_{US,t+1})n^*_{US,t}
+(Q_{W,t+1}+D_{W,t+1})n^*_{W,t}
+B^*_{US,t+1}+B^*_{W,t+1}.
\end{align*}
The term $\chi>0$ generates a convenience yield from holding US bonds and the ``exorbitant privilege'' of US bonds as a consequence.

\paragraph{Stock pricing (RoW effective SDF).}
The FOCs for equity positions imply, for $i\in\{US,W\}$,
\[
Q_{i,t}
=
E_t\!\left[\tilde M^{*,i}_{t,t+1}\bigl(Q_{i,t+1}+D_{i,t+1}\bigr)\right],
\]
with
\begin{align*}
\tilde M^{*,US}_{t,t+1}=\frac{M^*_{t,t+1}}{1+\phi_t^*(1-\omega_t^*)},
\qquad
\tilde M^{*,W}_{t,t+1}=\frac{M^*_{t,t+1}}{1-\phi_t^*\omega_t^*},\\
M^*_{t,t+1}
=\frac{\beta}{1-\beta}\,
\frac{C^*_{y,t}(C^*_{o,t+1})^{-\gamma}}{E_t\!\left[(C^*_{o,t+1})^{1-\gamma}\right]},
\qquad
\phi_t^* \equiv \frac{\kappa C^*_{y,t}}{(1-\beta)S_t^*}\,(\omega_t^*-\bar\omega^*).
\end{align*}

\paragraph{Saving notation and portfolio shares.}
Define present values and total saving:
\[
b_{US,t}^*\equiv q_{US,t}B^*_{US,t+1},
\qquad
b_{W,t}^*\equiv q_{W,t}B^*_{W,t+1},
\qquad
A_t^*\equiv S_t^*+b_{US,t}^*+b_{W,t}^*,
\qquad
C^*_{y,t}=e_{W,t}-A_t^*.
\]
Define bond shares in total saving:
\[
\theta_{US,t}^*\equiv \frac{b_{US,t}^*}{A_t^*},
\qquad
\theta_{W,t}^*\equiv \frac{b_{W,t}^*}{A_t^*},
\qquad
S_t^*=(1-\theta_{US,t}^*-\theta_{W,t}^*)A_t^*.
\]
Define the gross return on total saving:
\[
R_{A,t+1}^*
\equiv
(1-\theta_{US,t}^*-\theta_{W,t}^*)R_{p,t+1}^*
+\theta_{US,t}^* R_{f,t}
+\theta_{W,t}^* R_{f,t}^W,
\]
where
\[
R_{f,t}\equiv \frac{1}{q_{US,t}},
\qquad
R_{f,t}^W\equiv \frac{1}{q_{W,t}},
\qquad
R_{p,t+1}^*=\omega_t^*R_{US,t+1}+(1-\omega_t^*)R_{W,t+1}.
\]
Then $C^*_{o,t+1}=A_t^*R_{A,t+1}^*$.

\paragraph{RoW saving rule.}
Because $\chi\log(q_{US,t}B_{US,t+1}^*)=\chi\log(\theta_{US,t}^*A_t^*)$ contributes $\chi\log A_t^*$,
the saving choice separates and yields
\begin{align}\label{RoW_saving_rule_bond_cost}
\frac{1-\beta}{e_{W,t}-A_t^*}=\frac{\beta+\chi}{A_t^*}
\quad\Rightarrow\quad
A_t^*=\frac{\beta+\chi}{1+\chi}\,e_{W,t},
\qquad
C^*_{y,t}=\frac{1-\beta}{1+\chi}\,e_{W,t}.
\end{align}

\paragraph{RoW portfolio FOCs.}
Define the normalized RoW portfolio kernel
\[
\mathcal M_{t,t+1}^* \equiv \frac{(R_{A,t+1}^*)^{-\gamma}}{E_t\!\left[(R_{A,t+1}^*)^{1-\gamma}\right]}.
\]
The FOC for $\omega_t^*$ is
\begin{align}\label{RoW_omega_FOC_mod}
\beta\bigl(1-\theta_{US,t}^*-\theta_{W,t}^*\bigr)\,
E_t\!\left[\mathcal M_{t,t+1}^*\bigl(R_{US,t+1}-R_{W,t+1}\bigr)\right]
=
\kappa(\omega_t^*-\bar\omega^*).
\end{align}
The FOC for the RoW bond share $\theta_{W,t}^*$ is
\begin{align}\label{RoW_thetaW_FOC_mod}
\beta\,E_t\!\left[\mathcal M_{t,t+1}^*\bigl(R_{f,t}^W-R_{p,t+1}^*\bigr)\right]
=
0.
\end{align}
The FOC for the US bond share $\theta_{US,t}^*$ is
\begin{align}\label{RoW_thetaUS_FOC_mod}
\beta\,E_t\!\left[\mathcal M_{t,t+1}^*\bigl(R_{f,t}-R_{p,t+1}^*\bigr)\right]
+\frac{\chi}{\theta_{US,t}^*}
=
0.
\end{align}
And the FOC for the US bond amount $B_{US,t+1}^*$ delivers the exorbitant privilege wedge:
\[
q_{US,t}
=
E_t\!\left[M^*_{t,t+1}\right]
+\frac{\chi}{1+\chi}\,\frac{e_{W,t}}{B^*_{US,t+1}}.
\]

%--------------------------------------------------------
% Market clearing (two bonds with zero net supply)
%--------------------------------------------------------
\subsection{Market clearing}

\paragraph{World stock market clearing.}
\begin{align}\label{stock_clear_mod}
\begin{split}
n_{US,t}+n^*_{US,t}=1,\\
n_{W,t}+n^*_{W,t}=1.
\end{split}
\end{align}

\paragraph{Bond market clearing.}
Assume both one-period riskless bonds are in zero net supply. Since only RoW trades the RoW bond, bond market
clearing is
\begin{align}\label{bond_clear_mod_zerosupply}
B_{US,t+1}+B^*_{US,t+1}=0,
\qquad
B^*_{W,t+1}=0.
\end{align}
Equivalently, in present values,
\[
b_t+b_{US,t}^*=0,
\qquad
b_{W,t}^*=0.
\]

\paragraph{Stock market clearing in value form.}
Using portfolio weights,
\begin{align}\label{stock_weight_clear_mod}
\begin{split}
Q_{US,t} &= \omega_t S_t + \omega_t^* S_t^*,\\
Q_{W,t}  &= (1-\omega_t) S_t + (1-\omega_t^*) S_t^* .
\end{split}
\end{align}

\paragraph{Aggregate market-cap identity.}
Since $S_t=A_t-b_t$ and $S_t^*=A_t^*-b_{US,t}^*-b_{W,t}^*$, we have
\[
Q_{US,t}+Q_{W,t}=S_t+S_t^*=A_t+A_t^*-(b_t+b_{US,t}^*)-b_{W,t}^*.
\]
Using bond clearing \eqref{bond_clear_mod_zerosupply} and the saving rules
\eqref{US_saving_rule_bond_cost} and \eqref{RoW_saving_rule_bond_cost}, it follows that
\begin{align}\label{Qsum_identity_mod_zerosupply}
\boxed{
Q_{US,t}+Q_{W,t}
=
\beta e_{US,t}
+\frac{\beta+\chi}{1+\chi}\,e_{W,t}.
}
\end{align}

\paragraph{Goods market clearing.}
\begin{align}\label{goods_clear_mod}
C_{y,t}+C^*_{y,t}+C_{o,t}+C^*_{o,t}
=
e_{US,t}+e_{W,t}+D_{US,t}+D_{W,t}.
\end{align}



\section{Existence of bubbles}\label{sec:existence_of_bubbles}

In this section we reserve $\mathcal K_{t,t+1}$ for the undistorted normalized portfolio kernel,
\[
\mathcal K_{t,t+1}
=
\frac{R_{A,t+1}^{-\gamma}}{E_t\!\left[R_{A,t+1}^{1-\gamma}\right]},
\]
and we let $\mathcal M_{t,t+1}$ denote the \emph{effective} kernel that prices the US stock. By
\eqref{US_AP_bond_cost},
\begin{equation}\label{effective_US_kernel}
\mathcal M_{t,t+1}
\equiv
\widetilde M^{\,US}_{t,t+1}
=
\frac{\mathcal K_{t,t+1}}{\Psi_t},
\qquad
\Psi_t
\equiv
1-\lambda_t+\phi_t(1-\omega_t),
\end{equation}
where
\[
\phi_t=\frac{\kappa}{\beta(1-\theta_t)}(\omega_t-\bar\omega),
\qquad
\lambda_t=\frac{\eta\theta_t}{\beta}\,\Upsilon'(\theta_t).
\]

Using the saving rule \eqref{US_saving_rule_bond_cost} and the fact that $C_{o,t+1}=A_tR_{A,t+1}$,
the US \emph{undistorted} SDF simplifies to
\[
M_{t,t+1}
=
\frac{\beta}{1-\beta}\,
\frac{C_{y,t}C_{o,t+1}^{-\gamma}}{E_t\!\left[C_{o,t+1}^{1-\gamma}\right]}
=
\frac{R_{A,t+1}^{-\gamma}}{E_t\!\left[R_{A,t+1}^{1-\gamma}\right]}
=
\mathcal K_{t,t+1}.
\]
Hence the only difference between the stock-pricing kernel $\mathcal M_{t,t+1}$ and the undistorted kernel
$\mathcal K_{t,t+1}$ is the time-$t$ portfolio wedge $\Psi_t^{-1}$.

We henceforth restrict attention to equilibria such that
\begin{equation}\label{kernel_regular}
\Psi_t>0
\qquad\text{a.s. for all }t,
\end{equation}
so that the effective pricing kernel $\mathcal M_{t,t+1}$ is well defined and strictly positive.

Also note that because the RoW utility includes $\chi\log(q_{US,t}B^*_{US,t+1})$, feasibility requires
$B^*_{US,t+1}>0$. With zero net supply of US bonds,
$B_{US,t+1}+B^*_{US,t+1}=0$ implies $B_{US,t+1}<0$, hence
\begin{equation}\label{theta_nonpos}
\theta_t=\frac{b_t}{A_t}<0
\quad\Rightarrow\quad
1-\theta_t>1.
\end{equation}

Then we show that $\Psi_t>0$ are very likely to hold under mild parameter restrictions.
Because $0\le \omega_t\le 1$, we have $1-\omega_t\in[0,1]$ and
\[
(\omega_t-\bar\omega)(1-\omega_t)
\ge -\bar\omega(1-\omega_t)
\ge -\bar\omega.
\]
Therefore
\[
\phi_t(1-\omega_t)
=
\frac{\kappa}{\beta(1-\theta_t)}(\omega_t-\bar\omega)(1-\omega_t)
\ge
-\frac{\kappa\bar\omega}{\beta(1-\theta_t)}.
\]
We further assume the function form of issuance costs of bonds, for now, as
 $\Upsilon(\theta_t)=-\log(1-\theta_t)+\theta_t^2/2$. Then $\Psi_t >0$ if
\begin{equation} \label{suff_condition_pos_wedge}
	\beta > \frac{\eta \theta_t}{1-\theta_t}+\eta \theta_t^2 +\frac{\kappa\bar\omega}{1-\theta_t},	
\end{equation}
as
\[
\Psi_t > 1-\frac{\eta \theta_t}{\beta} \Upsilon'(\theta_t)-\frac{\kappa\bar\omega}{\beta(1-\theta_t)}.
\]
It is immediate that \eqref{suff_condition_pos_wedge} easily hold for sufficiently high $\beta$ and 
sufficiently low $\eta$ and $\kappa$ as $\theta_t$ usually falls within $[-1,0)$ and $1-\theta_t>1$.

\subsection{Temporary unbalanced growth and stochastic bubbles}

Following \citet{hirano2025technological}, let
\[
\mathcal{M}_{t,t+s}:=\prod_{j=0}^{s-1} \mathcal{M}_{t+j,t+j+1}.
\]
The US stock pricing equation \eqref{US_AP_bond_cost} (for $i=US$) implies
\[
Q_{US,0}=
\mathrm{E}_0 \left[\sum_{s=1}^{\infty} \mathcal{M}_{0,s} D_{US,s}\right]
+\lim_{t \rightarrow \infty} \mathrm{E}_0\left[\mathcal{M}_{0,t} Q_{US,t}\right],
\]
or more generally,
\[
Q_{US,t}=
\underbrace{\mathrm{E}_t \left[ \sum_{s=1}^{\infty} \mathcal{M}_{t,t+s} D_{US,t+s}\right]}_{\text{fundamental value }V_{US,t}}
+\underbrace{\lim_{T \rightarrow \infty} \mathrm{E}_t\left[\mathcal{M}_{t,T} Q_{US,T}\right]}_{\text{bubble }B_{US,t}}.
\]

Consider a regime-switching model with an absorbing state, as in \citet{weil1987confidence}.

\begin{ass}\label{ass_regime_switch}
There are two states of the US economy denoted by $u, b$. Letting $z_t \in \{u, b\}$ denote the state at time $t$,
the transition probabilities are given by
\[
P[z_{t+1} = u \mid z_t = u] = \pi \in (0,1), \qquad P[z_{t+1} = b \mid z_t = b] = 1.
\]
\end{ass}

\Cref{ass_regime_switch} implies that state $b$ is absorbing. We further assume that once state $b$ is reached, uncertainty is resolved,
and endowments and dividends grow at the same rate (balanced growth). On the other hand, the RoW is always on its balanced growth path (BGP hereafter).

\begin{ass}\label{ass_balanced_growth}
For any $\tau$, conditional on $z_\tau = b$, the sequence
$\{e_{US,t}, D_{US,t}, e_{W,t}, D_{W,t}\}_{t=\tau}^\infty$ is deterministic and
\[
\frac{e_{US,t+1}}{e_{US, t}} =
\frac{D_{US, t+1}}{D_{US,t}}=
\frac{e_{W,t+1}}{e_{W, t}} =
\frac{D_{W, t+1}}{D_{W,t}} ,\quad \forall t \geq \tau.
\]
\end{ass}

We can show that once state $b$ is reached, there is no bubble in the US stock market.

\begin{prop}\label{prop1}
If \Cref{ass_regime_switch,ass_balanced_growth} hold, once state $b$ is reached, the US stock price equals its fundamental value:
$Q_{US,t}=V_{US,t}$.
\end{prop}

\begin{proof}
The proof directly follows from \citet{hirano2025technological} because the US stock continues to satisfy a standard
Euler equation with the effective kernel $\mathcal M_{t,t+1}$, and state $b$ is deterministic balanced growth.
\end{proof}

By \Cref{prop1}, bubbles cannot arise in state $b$. So we assume $z_0=u$ from now on. Let $\tau$ be the stopping time that represents the first time the US economy switches to state $b$:
$\tau =\min \{t:z_t=b\}$. By \Cref{ass_regime_switch}, $\tau$ follows a geometric distribution,
\[
\mathrm{P}[\tau=j]=\pi^{j-1}(1-\pi), \quad j=1,2,\dots.
\]
Then the bubble term can be written as
\[
\mathrm{E}_0\left[\mathcal{M}_{0 , t} Q_{US,t}\right]
=\sum_{j=1}^t \pi^{j-1}(1-\pi) \mathrm{E}_0\left[\mathcal{M}_{0,t} Q_{US,t} \mid \tau=j\right]
+\pi^t \mathrm{E}_0\left[\mathcal{M}_{0, t} Q_{US,t} \mid \tau>t\right].
\]

The following proposition provides a sufficient condition for the nonexistence of bubbles.

\begin{prop}\label{prop2}
Suppose \Cref{ass_regime_switch,ass_balanced_growth} hold. Then there is no bubble if and only if
\[
\lim _{t \rightarrow \infty} \pi^t \mathrm{E}_0\left[\mathcal{M}_{0,t} Q_{US,t} \mid \tau>t\right]=0.
\]
If $(e_{US,t}, D_{US,t})$ is known at time $t-1$ (i.e., $(e_{US,t},D_{US,t})$ is $\mathcal{F}_{t-1}$-measurable), then there is no bubble.
\end{prop}

\begin{proof}
The proof directly follows from \citet{hirano2025technological}.
\end{proof}

Aligned with the sufficient conditions in \Cref{prop2}, we introduce the following assumption.

\begin{ass}\label{ass_switch_dependence}
Conditional on time $t-1$ information, the US endowment $e_{US,t}$ and US dividend $D_{US,t}$ depend only on the state
$z_t \in \{u,b\}$.
\end{ass}


Now we are ready to present a more general but less primitive version of Theorem 1 in \citet{hirano2025technological}.
For $z\in\{u,b\}$, let $x_t^z$ be the value of $x_{US,t}^z$ conditional on $z_0=\cdots=z_{t-1}=u$ and $z_t=z$,
and let $C_t^z=C_{o,t}^z$ denote US old-age consumption at date $t$ on branch $z$.

\begin{thm}\label{thm_1}
Suppose Assumptions \ref{ass_regime_switch}, \ref{ass_balanced_growth}, and \ref{ass_switch_dependence} hold. If $z_0=u$, then there is a bubble at $t=0$ if and only if
\begin{subequations}\label{thm_cond1}
\begin{align}
\sum_{t=1}^\infty \frac{D_t^u}{Q_t^u}&<\infty, \label{thm_cond_1a}\\
\sum_{t=1}^\infty
\left(\frac{C_t^b}{C_t^u}\right)^{1-\gamma}\;
\frac{Q_t^b+D_t^b}{C_t^b}\;
\frac{C_t^u}{Q_t^u}&<\infty. \label{thm_cond_1b}
\end{align}
\end{subequations}
\end{thm}

\begin{proof}
Let
\[
B_t \equiv \mathbb E_0\!\left[\mathcal M_{0,t}\,Q_t^u\,\mathbf 1\{\tau>t\}\right].
\]
Because $\mathbb P(\tau>t)=\pi^t$, this is equivalent to
\[
B_t=\pi^t\,\mathbb E_0\!\left[\mathcal M_{0,t}\,Q_t^u\mid \tau>t\right].
\]
By \Cref{prop2}, there is a bubble at $0$ if and only if $\lim_{t\to\infty}B_t>0$ (and no bubble if and only if $\lim_{t\to\infty}B_t=0$).

On histories with $z_{t-1}=u$, the US stock Euler equation is
\[
Q_{t-1}^u
=\mathbb E_{t-1}\!\left[\mathcal M_{t-1,t}\,(Q_t^{z}+D_t^{z})\right],
\qquad z\in\{u,b\}.
\]
Multiply both sides by $\mathcal M_{0,t-1}\mathbf 1\{\tau>t-1\}$ and take $\mathbb E_0$. Using
$\mathcal M_{0,t}=\mathcal M_{0,t-1}\mathcal M_{t-1,t}$, we obtain the exact identity
\begin{equation}\label{recursion_bubble_1}
\mathbb E_0\!\left[\mathcal M_{0,t-1}Q_{t-1}^u\mathbf 1\{\tau>t-1\}\right]
=
\mathbb E_0\!\left[\mathcal M_{0,t}(Q_t^u+D_t^u)\mathbf 1\{\tau>t\}\right]
+
\mathbb E_0\!\left[\mathcal M_{0,t}(Q_t^b+D_t^b)\mathbf 1\{\tau=t\}\right].
\end{equation}
Split the first term on the right-hand side to get
\[
\mathbb E_0\!\left[\mathcal M_{0,t}(Q_t^u+D_t^u)\mathbf 1\{\tau>t\}\right]
=
B_t+\mathbb E_0\!\left[\mathcal M_{0,t}D_t^u\mathbf 1\{\tau>t\}\right].
\]
Thus \eqref{recursion_bubble_1} becomes
\begin{equation}\label{bubble_recursion_2}
B_{t-1}
=
B_t
+\underbrace{\mathbb E_0\!\left[\mathcal M_{0,t}D_t^u\mathbf 1\{\tau>t\}\right]}_{\text{dividend leakage term}}
+\underbrace{\mathbb E_0\!\left[\mathcal M_{0,t}(Q_t^b+D_t^b)\mathbf 1\{\tau=t\}\right]}_{
	\text{switch-at-t leakage term}},
\end{equation}
where the last two terms can be interpreted as leakage terms that prevent the bubble from growing too fast. 
In other words, \textit{leakage} is defined as the amount of present value that drains out of the 
continuation bubble mass when we move one period forward. It’s exactly the gap $B_{t-1}-B_t$.

Define the leakage ratio
\[
a_t\equiv
\frac{
\mathbb E_0\!\left[\mathcal M_{0,t}D_t^u\mathbf 1\{\tau>t\}\right]
+\mathbb E_0\!\left[\mathcal M_{0,t}(Q_t^b+D_t^b)\mathbf 1\{\tau=t\}\right]
}{
\mathbb E_0\!\left[\mathcal M_{0,t}Q_t^u\mathbf 1\{\tau>t\}\right]
}
=\frac{\text{two leakage terms}}{B_t}\ge 0.
\]
Then \eqref{bubble_recursion_2} implies $B_{t-1}=(1+a_t)B_t$, so iterating yields
\[
B_t=B_0\prod_{s=1}^t\frac{1}{1+a_s},
\qquad B_0=Q_0^u>0.
\]
Since $a_s\ge 0$,
\[
B_0\left(1+\sum_{s=1}^t a_s\right)^{-1}
\ge B_0\prod_{s=1}^t\frac{1}{1+a_s}
\ge B_0\exp\!\left(-\sum_{s=1}^t a_s\right).
\]
Hence,
\[
\lim_{t\to\infty}B_t>0
\quad\Longleftrightarrow\quad
\sum_{t=1}^\infty a_t<\infty.
\]

It remains to express $a_t$ in terms of observables along the $u$-history and at the switching date.
On $\{\tau>t\}$ we have $z_1=\cdots=z_t=u$, hence
\[
\mathbb E_0\!\left[\mathcal M_{0,t}D_t^u\mathbf 1\{\tau>t\}\right]
=\pi^t\,\mathcal M_{0,t}^u\,D_t^u,
\qquad
B_t=\pi^t\,\mathcal M_{0,t}^u\,Q_t^u,
\]
so
\[
\frac{\mathbb E_0[\mathcal M_{0,t}D_t^u\,\mathbf 1\{\tau>t\}]}{B_t}
=\frac{D_t^u}{Q_t^u}.
\]

On $\{\tau=t\}$, we have $z_1=\cdots=z_{t-1}=u$ and $z_t=b$. Hence
\[
\mathbb E_0\!\left[\mathcal M_{0,t}(Q_t^b+D_t^b)\mathbf 1\{\tau=t\}\right]
=\pi^{t-1}(1-\pi)\,\mathcal M_{0,t-1}^u\,\mathcal M_{t-1,t}^b\,(Q_t^b+D_t^b).
\]
Dividing by $B_t=\pi^{t}\mathcal M_{0,t-1}^u\mathcal M_{t-1,t}^u Q_t^u$ gives
\[
\frac{\mathbb E_0[\mathcal M_{0,t}(Q_t^b+D_t^b)\mathbf 1\{\tau=t\}]}{B_t}
=
\frac{1-\pi}{\pi}\;
\frac{\mathcal M_{t-1,t}^b}{\mathcal M_{t-1,t}^u}\;
\frac{Q_t^b+D_t^b}{Q_t^u}.
\]

Finally, under \Cref{ass_switch_dependence}, the portfolio choices $(\omega_{t-1},\theta_{t-1})$ are $\mathcal F_{t-1}$-measurable and do not depend on $z_t$.
Thus the full time-$(t-1)$ denominator
\[
\Psi_{t-1}=1-\lambda_{t-1}+\phi_{t-1}(1-\omega_{t-1})
\]
is $\mathcal F_{t-1}$-measurable and common across $z_t\in\{u,b\}$, so it cancels in the ratio:
\[
\frac{\mathcal M_{t-1,t}^b}{\mathcal M_{t-1,t}^u}
=
\frac{\mathcal K_{t-1,t}^b}{\mathcal K_{t-1,t}^u}
=
\frac{(C_t^b)^{-\gamma}}{(C_t^u)^{-\gamma}}
=
\left(\frac{C_t^u}{C_t^b}\right)^{\gamma}.
\]
Therefore,
\[
a_t
=
\frac{D_t^u}{Q_t^u}
+\frac{1-\pi}{\pi}\left(\frac{C_t^u}{C_t^b}\right)^{\gamma}\frac{Q_t^b+D_t^b}{Q_t^u}
=
\frac{D_t^u}{Q_t^u}
+\frac{1-\pi}{\pi}
\left(\frac{C_t^b}{C_t^u}\right)^{1-\gamma}
\frac{Q_t^b+D_t^b}{C_t^b}\frac{C_t^u}{Q_t^u}.
\]
Hence $\sum_{t=1}^\infty a_t<\infty$ is equivalent to \eqref{thm_cond1}.
\end{proof}

Intuitively, the first summability condition \eqref{thm_cond_1a} requires that dividends along the $u$-history 
do not grow too fast relative to stock prices, while the second summability condition \eqref{thm_cond_1b} 
requires that the discounted value of dividends and stock prices at the switching date does not grow too 
fast relative to old-age consumption along the $u$-history. Therefore, the bubble candidate survives only if,
 over time, dividends and switching events do not drain too much value relative to what remains; that is, 
 only if the cumulative leakage $\sum a_t$ is finite.

\subsection{Primitive sufficient and necessary conditions}

We now derive primitive conditions for the existence of bubbles by translating the two summability conditions
in \Cref{thm_1} into exogenous variables along the $u$-history and at the switching date. We first impose primitive
boundedness restrictions on relative country size, switch-date US fundamentals, and dividends.

\begin{ass}[Relative country size and dividend bounds]\label{ass_relative_dividend_bounds}
There exists $\overline\lambda_{e_W}<\infty$ and $\overline\lambda_{D_W}<\infty$ such that
\[
\frac{e_{W,t}}{e_{US,t}^u}\le \overline\lambda_{e_W},
\qquad 
\frac{D_{W,t}}{e_{US,t}^u} \le \overline\lambda_{D_W},
\qquad\text{for all }t\le \tau.
\]
\end{ass}

\begin{ass}[Switch-date US fundamentals]\label{ass_switch_us_fundamentals}
There exist constants $0<\underline\lambda_e\le \overline\lambda_e<\infty$ and
$0\le \overline\lambda_{D_{US}}<\infty$ such that, whenever $z_{t-1}=u$ and $z_t=b$,
\[
\underline\lambda_e
\le
\frac{e_{US,t}^b}{e_{US,t}^u}
\le
\overline\lambda_e,
\qquad
\frac{D_{US,t}^b}{e_{US,t}^u}
\le
\overline\lambda_{D_{US}}.
\]
\end{ass}


Finally, as shown in \eqref{theta_nonpos}, the convenience-yield term forces the US to be a net borrower:
$\theta_t\le 0$ and $1-\theta_t\ge 1$. With these ingredients, we can state a single corollary that 
gives primitive conditions ensuring that the
summability conditions in \Cref{thm_1} reduce to \eqref{thm_cond2}.

\begin{coro}\label{coro1}
Suppose \Crefrange{ass_issuance_costs}{ass_switch_us_fundamentals} hold. If:

(i) equity weights are uniformly interior along both relevant branches, i.e., there exists
$\underline{\lambda}_\omega>0$ such that
\[
\omega_t^u\ge \underline{\lambda}_\omega,
\qquad
\omega_t^b\ge \underline{\lambda}_\omega,
\qquad\text{for all }t\le\tau;
\]

(ii) the bond-share FOC \eqref{US_theta_FOC_bond_cost} holds with $\eta>0$ and finite LHS, i.e.,
\[
\left|\,\beta\,E_t\!\left[\mathcal K_{t,t+1}\bigl(R_{f,t}-R_{p,t+1}\bigr)\right]\right|
\le \overline\xi;
\]

(iii) $R_{f,t}$ is uniformly bounded above on $t\le\tau$, i.e.,
\[
R_{f,t}\le \overline R_f<\infty;
\]

(iv) old-age consumption is uniformly comparable to $e_{US,t}^u$, i.e., there exist constants
$0<\underline{\lambda}_C\le \overline{\lambda}_C<\infty$ such that
\[
\underline{\lambda}_C\,e_{US,t}^u
\le
C_t^z
\le
\overline{\lambda}_C\,e_{US,t}^u,
\qquad
z\in\{u,b\},\ \forall t\le\tau,
\]

then the two conditions \eqref{thm_cond1} in \Cref{thm_1} are equivalent to the primitive conditions
\begin{subequations}\label{thm_cond2}
\begin{align}
\sum_{t=1}^\infty \frac{D_t^u}{e_{US,t}^u}&<\infty, \label{thm_cond_2a}\\
\sum_{t=1}^\infty
\left(\frac{C_t^b}{C_t^u}\right)^{1-\gamma}&<\infty. \label{thm_cond_2b}
\end{align}
\end{subequations}
\end{coro}

\begin{proof}
We show (a) \eqref{thm_cond_1a} $\Longleftrightarrow$ \eqref{thm_cond_2a} and
(b) \eqref{thm_cond_1b} $\Longleftrightarrow$ \eqref{thm_cond_2b}.

\medskip
\noindent\textbf{(a) Proof of \eqref{thm_cond_1a} $\Longleftrightarrow$ \eqref{thm_cond_2a}.}
By \eqref{Qsum_identity_mod_zerosupply} and \Cref{ass_relative_dividend_bounds},
\[
Q_{US,t}^u \le Q_{US,t}^u+Q_{W,t}^u
=\beta e_{US,t}^u+\frac{\beta+\chi}{1+\chi}e_{W,t}
\le \left(\beta+\frac{\beta+\chi}{1+\chi}\overline\lambda_{e_W}\right)e_{US,t}^u,
\qquad t\le\tau.
\]
Moreover, by \eqref{stock_weight_clear_mod},
\[
Q_{US,t}^u=\omega_t^u S_t^u+\omega_t^{*,u}S_t^{*,u}\ge \omega_t^u S_t^u.
\]
Using $S_t^u=(1-\theta_t^u)A_t^u=(1-\theta_t^u)\beta e_{US,t}^u$ and \eqref{theta_nonpos}, we obtain
\[
Q_{US,t}^u\ge \underline\lambda_\omega\,\beta\,e_{US,t}^u
\qquad\text{for all }t\le\tau.
\]
Hence there exist constants $0<\underline q\le \bar q<\infty$ such that
\[
\underline q\,e_{US,t}^u\le Q_{US,t}^u\le \bar q\,e_{US,t}^u
\qquad\text{for all }t\le\tau,
\]
which implies
\[
\sum_{t=1}^\infty \frac{D_t^u}{Q_t^u}<\infty
\quad\Longleftrightarrow\quad
\sum_{t=1}^\infty \frac{D_t^u}{e_{US,t}^u}<\infty.
\]

\medskip
\noindent\textbf{(b) Proof of \eqref{thm_cond_1b} $\Longleftrightarrow$ \eqref{thm_cond_2b}.}
Define
\[
m_t:=\frac{Q_t^b+D_t^b}{C_t^b}\cdot \frac{C_t^u}{Q_t^u}.
\]
It suffices to show that $m_t$ is bounded above and below by strictly positive constants uniformly over $t\le\tau$.

From part (a), there exist constants $0<\underline q\le \overline q<\infty$ such that
\[
\underline q\,e_{US,t}^u\le Q_t^u\le \overline q\,e_{US,t}^u
\qquad\text{for all }t\le\tau.
\]
Next, by \eqref{stock_weight_clear_mod}, condition (i), the US saving rule, \eqref{theta_nonpos}, and \Cref{ass_switch_us_fundamentals},
\[
Q_t^b=\omega_t^b S_t^b+\omega_t^{*,b}S_t^{*,b}
\ge \omega_t^b S_t^b
\ge \underline\lambda_\omega (1-\theta_t^b)A_t^b
\ge \underline\lambda_\omega\,\beta\,e_{US,t}^b
\ge \underline\lambda_\omega\,\beta\,\underline\lambda_e\,e_{US,t}^u.
\]
Also, by \eqref{Qsum_identity_mod_zerosupply} and \Cref{ass_relative_dividend_bounds,ass_switch_us_fundamentals},
\[
Q_t^b\le Q_t^b+Q_{W,t}^b
=\beta e_{US,t}^b+\frac{\beta+\chi}{1+\chi}e_{W,t}
\le
\left(\beta\overline\lambda_e+\frac{\beta+\chi}{1+\chi}\overline\lambda_{e_W}\right)e_{US,t}^u.
\]
Using again \Cref{ass_switch_us_fundamentals}, we obtain
\[
\underline\lambda_\omega\,\beta\,\underline\lambda_e\,e_{US,t}^u
\le Q_t^b+D_t^b
\le
\left(\beta\overline\lambda_e+\frac{\beta+\chi}{1+\chi}\overline\lambda_{e_W}+\overline\lambda_{D_{US}}\right)e_{US,t}^u.
\]
Therefore there exist constants $0<\underline q_b\le \overline q_b<\infty$ such that
\[
\underline q_b\,e_{US,t}^u
\le
Q_t^b+D_t^b
\le
\overline q_b\,e_{US,t}^u
\qquad\text{for all }t\le\tau.
\]

By condition (iv),
\[
\underline\lambda_C\,e_{US,t}^u
\le
C_t^z
\le
\overline\lambda_C\,e_{US,t}^u,
\qquad z\in\{u,b\},\ \forall t\le\tau.
\]
Hence
\[
\frac{\underline q_b}{\overline\lambda_C}
\le
\frac{Q_t^b+D_t^b}{C_t^b}
\le
\frac{\overline q_b}{\underline\lambda_C},
\qquad
\frac{\underline\lambda_C}{\overline q}
\le
\frac{C_t^u}{Q_t^u}
\le
\frac{\overline\lambda_C}{\underline q}.
\]
Thus $m_t$ is bounded above and below by strictly positive constants uniformly over $t\le\tau$.
Therefore,
\[
\sum_{t=1}^\infty
\left(\frac{C_t^b}{C_t^u}\right)^{1-\gamma}
\frac{Q_t^b+D_t^b}{C_t^b}
\frac{C_t^u}{Q_t^u}
<\infty
\quad\Longleftrightarrow\quad
\sum_{t=1}^\infty
\left(\frac{C_t^b}{C_t^u}\right)^{1-\gamma}
<\infty.
\]
This proves \eqref{thm_cond_1b} $\Longleftrightarrow$ \eqref{thm_cond_2b}.
\end{proof}



%========================================================
% Appendix
%========================================================

\appendix
%\section{Proof of Lemma~\ref{lem1}}


\bibliography{references.bib} 

\end{document}